% GENERATED FILE. Edit canonical lessons, then run npm run generate:books.
% canonical-source-hash: fnv1a64:4dca3509c3afd520
% canonical-lessons: AR-C469-turba, AR-C469-samad, AR-C469-mahsul, AR-C469-khasib, AR-C469-iqtalaa

\chapter{Soil and crops}
\label{ch:ar-soil-and-crops}
\hlchaptermodality{469}

\begin{chapteropening}
\textbf{By the end of this chapter:} \emph{I can say soil, fertiliser, a crop, fertile and to uproot.}

Complete the last lesson of chapter 469: I can say soil, fertiliser, a crop, fertile and to uproot.
\end{chapteropening}

\section[\texorpdfstring{\ar{تربة} (turba)}{تربة (turba)}]{\ar{تربة} (turba) — soil; earth}

\label{lesson:AR-C469-turba}



\begin{quote}
Before the new one: say the Arabic for to bloom, then the Arabic for to bear fruit.
\end{quote}

\subsection*{You'll want to know: \ar{تربة}}
\textbf{\ar{تربة}} — \emph{turba} — \textquotedblleft{}soil; earth\textquotedblright{}.

\begin{quote}
\ar{التربة} \ar{هنا} \ar{جيدة} \ar{للقمح}.
\end{quote}

\begin{quote}
\emph{at-turba hunā jayyida li-l-qamḥ.}
\end{quote}

\begin{quote}
The soil here is good for wheat.
\end{quote}

\begin{cousinweb}
\textbf{\ar{تربة}} (\emph{turba}) is built on the native Arabic root \emph{t-r-b}, \textquotedblleft{}earth, dust\textquotedblright{} (Lane). It shares its root with \emph{turāb}, \textquotedblleft{}dust, soil\textquotedblright{}.
\end{cousinweb}

\begin{grammarlens}[title={What you've built}]
One more word for getting around.
\end{grammarlens}

\subsection*{Your turn}
\begin{itemize}
\raggedright
  \item \emph{Say it:} \emph{turba}
  \item \emph{Say it:} \emph{turba}, once more
  \item \emph{Say it:} \emph{at-turba hunā jayyida li-l-qamḥ.} — The soil here is good for wheat.
  \item \emph{Recall:} say the Arabic for to bloom, then the Arabic for to bear fruit, then say \emph{turba} again
\end{itemize}

\begin{tcolorbox}[breakable,colback=teal!4,colframe=teal!35!black,title={Before you move on}]
What does \ar{تربة} mean? (\textquotedblleft{}Soil; earth\textquotedblright{}.) Say it once more.
\end{tcolorbox}
\section[\texorpdfstring{\ar{سماد} (samād)}{سماد (samād)}]{\ar{سماد} (samād) — fertiliser; manure}

\label{lesson:AR-C469-samad}



\begin{quote}
Before the new one: say the Arabic for to bear fruit, then the Arabic for soil.
\end{quote}

\subsection*{You'll want to know: \ar{سماد}}
\textbf{\ar{سماد}} — \emph{samād} — \textquotedblleft{}fertiliser; manure\textquotedblright{}.

\begin{quote}
\ar{وضع} \ar{المزارع} \ar{السماد} \ar{في} \ar{الحقل}.
\end{quote}

\begin{quote}
\emph{waḍaʿa al-muzāriʿ as-samād fī al-ḥaql.}
\end{quote}

\begin{quote}
The farmer put fertiliser on the field.
\end{quote}

\begin{cousinweb}
\textbf{\ar{سماد}} (\emph{samād}) is built on the native Arabic root \emph{s-m-d}, \textquotedblleft{}to manure the land\textquotedblright{} (Lane). Lane gives \emph{samād} as manure of dung and earth.
\end{cousinweb}

\begin{grammarlens}[title={What you've built}]
One more word for getting around.
\end{grammarlens}

\subsection*{Your turn}
\begin{itemize}
\raggedright
  \item \emph{Say it:} \emph{samād}
  \item \emph{Say it:} \emph{samād}, once more
  \item \emph{Say it:} \emph{waḍaʿa al-muzāriʿ as-samād fī al-ḥaql.} — The farmer put fertiliser on the field.
  \item \emph{Recall:} say the Arabic for to bear fruit, then the Arabic for soil, then say \emph{samād} again
\end{itemize}

\begin{tcolorbox}[breakable,colback=teal!4,colframe=teal!35!black,title={Before you move on}]
What does \ar{سماد} mean? (\textquotedblleft{}Fertiliser; manure\textquotedblright{}.) Say it once more.
\end{tcolorbox}
\section[\texorpdfstring{\ar{محصول} (maḥṣūl)}{محصول (maḥṣūl)}]{\ar{محصول} (maḥṣūl) — a crop; a harvest}

\label{lesson:AR-C469-mahsul}



\begin{quote}
Before the new one: say the Arabic for soil, then the Arabic for fertiliser.
\end{quote}

\subsection*{You'll want to know: \ar{محصول}}
\textbf{\ar{محصول}} — \emph{maḥṣūl} — \textquotedblleft{}a crop; a harvest\textquotedblright{}.

\begin{quote}
\ar{كان} \ar{محصول} \ar{القمح} \ar{كبيرا}.
\end{quote}

\begin{quote}
\emph{kāna maḥṣūl al-qamḥ kabīran.}
\end{quote}

\begin{quote}
The wheat crop was large.
\end{quote}

Plural: \ar{محاصيل}, \emph{maḥāṣīl}.

\begin{cousinweb}
\textbf{\ar{محصول}} (\emph{maḥṣūl}) is built on the native Arabic root \emph{ḥ-ṣ-l}, \textquotedblleft{}to remain, to be obtained\textquotedblright{} (Lane). It is a passive participle: what is obtained.
\end{cousinweb}

\begin{grammarlens}[title={What you've built}]
One more word for getting around.
\end{grammarlens}

\subsection*{Your turn}
\begin{itemize}
\raggedright
  \item \emph{Say it:} \emph{maḥṣūl}
  \item \emph{Say it:} \emph{maḥṣūl}, once more
  \item \emph{Say it:} \emph{kāna maḥṣūl al-qamḥ kabīran.} — The wheat crop was large.
  \item \emph{Recall:} say the Arabic for soil, then the Arabic for fertiliser, then say \emph{maḥṣūl} again
\end{itemize}

\begin{tcolorbox}[breakable,colback=teal!4,colframe=teal!35!black,title={Before you move on}]
What does \ar{محصول} mean? (\textquotedblleft{}A crop; a harvest\textquotedblright{}.) Say it once more.
\end{tcolorbox}
\section[\texorpdfstring{\ar{خصب} (khaṣib)}{خصب (khaṣib)}]{\ar{خصب} (khaṣib) — fertile}

\label{lesson:AR-C469-khasib}



\begin{quote}
Before the new one: say the Arabic for fertiliser, then the Arabic for a crop.
\end{quote}

\subsection*{You'll want to know: \ar{خصب}}
\textbf{\ar{خصب}} — \emph{khaṣib} — \textquotedblleft{}fertile\textquotedblright{}.

\begin{quote}
\ar{الأرض} \ar{حول} \ar{النهر} \ar{خصبة}.
\end{quote}

\begin{quote}
\emph{al-arḍ ḥawla an-nahr khaṣiba.}
\end{quote}

\begin{quote}
The land around the river is fertile.
\end{quote}

The noun \emph{khiṣb}, \textquotedblleft{}fertility\textquotedblright{}, has the same letters.

\begin{cousinweb}
\textbf{\ar{خصب}} (\emph{khaṣib}) is built on the native Arabic root \emph{kh-ṣ-b}, \textquotedblleft{}to be fruitful, rich in herbage\textquotedblright{} (Lane).
\end{cousinweb}

\begin{grammarlens}[title={What you've built}]
One more word for getting around.
\end{grammarlens}

\subsection*{Your turn}
\begin{itemize}
\raggedright
  \item \emph{Say it:} \emph{khaṣib}
  \item \emph{Say it:} \emph{khaṣib}, once more
  \item \emph{Say it:} \emph{al-arḍ ḥawla an-nahr khaṣiba.} — The land around the river is fertile.
  \item \emph{Recall:} say the Arabic for fertiliser, then the Arabic for a crop, then say \emph{khaṣib} again
\end{itemize}

\begin{tcolorbox}[breakable,colback=teal!4,colframe=teal!35!black,title={Before you move on}]
What does \ar{خصب} mean? (\textquotedblleft{}Fertile\textquotedblright{}.) Say it once more.
\end{tcolorbox}
\section[\texorpdfstring{\ar{اقتلع} (iqtalaʿa)}{اقتلع (iqtalaʿa)}]{\ar{اقتلع} (iqtalaʿa) — to uproot}

\label{lesson:AR-C469-iqtalaa}



\begin{quote}
Before the new one: say the Arabic for a crop, then the Arabic for fertile.
\end{quote}

\subsection*{You'll want to know: \ar{اقتلع}}
\textbf{\ar{اقتلع}} — \emph{iqtalaʿa} — \textquotedblleft{}to uproot\textquotedblright{}.

\begin{quote}
\ar{اقتلعت} \ar{الريح} \ar{شجرة} \ar{كبيرة}.
\end{quote}

\begin{quote}
\emph{iqtalaʿat ar-rīḥ shajara kabīra.}
\end{quote}

\begin{quote}
The wind uprooted a big tree.
\end{quote}

Present: \ar{يقتلع}, \emph{yaqtaliʿ}.

\begin{cousinweb}
\textbf{\ar{اقتلع}} (\emph{iqtalaʿa}) is built on the native Arabic root \emph{q-l-ʿ}, \textquotedblleft{}to pull up, to uproot\textquotedblright{} (Lane). It is Form VIII.
\end{cousinweb}

\begin{grammarlens}[title={What you've built}]
One more word for getting around.
\end{grammarlens}

\subsection*{Your turn}
\begin{itemize}
\raggedright
  \item \emph{Say it:} \emph{iqtalaʿa}
  \item \emph{Say it:} \emph{iqtalaʿa}, once more
  \item \emph{Say it:} \emph{iqtalaʿat ar-rīḥ shajara kabīra.} — The wind uprooted a big tree.
  \item \emph{Recall:} say the Arabic for a crop, then the Arabic for fertile, then say \emph{iqtalaʿa} again
\end{itemize}

\begin{tcolorbox}[breakable,colback=teal!4,colframe=teal!35!black,title={Before you move on}]
What does \ar{اقتلع} mean? (\textquotedblleft{}To uproot\textquotedblright{}.) Say it once more.
\end{tcolorbox}
